\chapter{预处理与宏}
\label{ch:19}

本章介绍以 \texttt{\%} 和 \texttt{\$} 开头的两类源码加工工具：Magel 的预处理和卫生宏。预处理控制源码是否参与编译，宏由解析器建立 AST 并在语义分析期间展开。

这一章分两半：

\begin{itemize}

\item 预处理（\texttt{\%} 开头）：条件编译——「这段代码要不要编进去」；

\item 宏（\texttt{\$} 开头）：卫生 AST 替换——保留绑定、类型与来源的编译期代码。

\end{itemize}

先说一句要紧的话：它们不污染普通源码。 预处理指令以 \texttt{\%} 开头、宏调用以 \texttt{\$} 开头，和普通标识符不共享命名空间——不会像 C 的 \texttt{\#define} 那样，一个叫 \texttt{MIN} 的宏把别人的变量名、结构体字段一起换掉。所以用不用只看一条：这段文本是不是得在编译前就改。 是，就放心用；不是，函数、泛型、常量通常更省事。

\section{条件编译：\texttt{\%define} + \texttt{\%if}}

同一个程序，调试时想多打日志、发布时想去掉——这是条件编译最典型的用途：

\par
\medskip
\begin{lstlisting}[style=magel]
%define DEBUG 1

int64 main() {
%if DEBUG
    println("debug on");
%else
    println("debug off");
%endif
    return 0;
}
\end{lstlisting}
\par
\medskip

\par
\medskip
\begin{lstlisting}[style=out]
debug on
\end{lstlisting}
\par
\medskip

把第一行改成 \texttt{\%define DEBUG 0}，同一份代码就输出另一个分支：

\par
\medskip
\begin{lstlisting}[style=magel]
%define DEBUG 0

int64 main() {
%if DEBUG
    println("debug on");
%else
    println("debug off");
%endif
    return 0;
}
\end{lstlisting}
\par
\medskip

\par
\medskip
\begin{lstlisting}[style=out]
debug off
\end{lstlisting}
\par
\medskip

几条要点：

\begin{itemize}

\item \texttt{\%define} 定义预处理符号；

\item \texttt{\%if 符号} 看它的值：非 0 走 \texttt{\%if} 分支，0 走 \texttt{\%else}；

\item 不活跃的分支整块丢弃——一个字都不会编进去；

\item 指令行必须顶格写在行首（\texttt{\%} 前面不能有别的字符）；

\item 被丢弃/指令占用的行用空行占位，所以预处理后行号和源文件一一对应，报错永远指回你真正写错的那一行。

\end{itemize}

\begin{tip}{提示}
\texttt{\%if} 后面要写符号名，不能写 \texttt{\%if 1} 这种字面量——写了会报「这里要写一个符号名」。这是它和 C 的 \texttt{\#if} 最容易混的地方。
\end{tip}

\section{按「有没有定义」分支：\texttt{\%ifdef} / \texttt{\%ifndef}}

有时候你只关心「这个符号定义过没有」，不关心值：

\par
\medskip
\begin{lstlisting}[style=magel]
int64 main() {
%ifdef FEATURE
    println("有 FEATURE");
%else
    println("没有 FEATURE");
%endif
    return 0;
}
\end{lstlisting}
\par
\medskip

\par
\medskip
\begin{lstlisting}[style=out]
没有 FEATURE
\end{lstlisting}
\par
\medskip

\begin{itemize}

\item \texttt{\%ifdef 符号}：定义过就走；

\item \texttt{\%ifndef 符号}：没定义才走（\texttt{\%ifdef} 的反面）。

\end{itemize}

\texttt{\%if} 和 \texttt{\%ifdef} 的分工记住一句话：\texttt{\%if} 看值，\texttt{\%ifdef} 看有没有。

\section{从命令行传开关：\texttt{-D}}

开关真正好用，是构建时才决定——同一份源码，编出不同版本：

\par
\medskip
\begin{lstlisting}[style=sh]
magel run app.mgl -D FEATURE
\end{lstlisting}
\par
\medskip

\par
\medskip
\begin{lstlisting}[style=out]
有 FEATURE
\end{lstlisting}
\par
\medskip

\texttt{-D FEATURE} 等价于在文件开头写了 \texttt{\%define FEATURE}。于是同一份代码：

\begin{itemize}

\item \texttt{magel run app.\allowbreak{}mgl} \ensuremath{\rightarrow} 「没有 FEATURE」；

\item \texttt{magel run app.\allowbreak{}mgl -D FEATURE} \ensuremath{\rightarrow} 「有 FEATURE」。

\end{itemize}

发布版 / 调试版、不同硬件平台，就靠这个开关切。

\section{条件编译的边界}

\texttt{\%define} 管理当前编译单元内的预处理符号，命令行 \texttt{-D} 提供各模块编译时使用的开关。不活跃分支不参与后续语言编译，但行号仍保留。

模块复用通过 \texttt{module} 和 \texttt{import}，完整规则归入\hyperref[ch:18]{第 18 章}。条件编译可控制导入语句是否参与编译，却不是把被导入文件当作文本贴到当前位置。

\section{卫生 AST 宏}

\texttt{\$macro} 定义由生成解析器建立语法树，语义分析按依赖展开，类型推导与展开交错进行，四元式生成前完成。宏不把替换结果拼成源码重新解析。参数保留原文、位置、语法候选及调用处绑定；引入的普通局部声明自动卫生，辅助名称在定义模块查找。

\par
\medskip
\begin{lstlisting}[style=magel]
$macro SQUARE(x)
x * x
$end_macro
$macro ANSWER
42
$end_macro
$macro EMPTY()
42
$end_macro
int64 main() {
    int64 temporary = 10;
    println("{} {} {}", $SQUARE(1 + 2), $ANSWER, $EMPTY());
    return 0;
}
\end{lstlisting}
\par
\medskip

实参 \texttt{1 + 2} 是完整表达式，替换进乘法的两个操作数，结果为 9。无需形参括号保护，也不自动给输出加括号、缓存参数或插入作用域；\texttt{\$SQUARE(\allowbreak{}next(\allowbreak{})\allowbreak{})\allowbreak{}} 仍调用两次 \texttt{next(\allowbreak{})\allowbreak{}}。编号临时变量、宏头 \texttt{expr} 标注及形参保护标注均已删除。语句宏的调用写完整分号；宏体的语句直接插入调用所在列表。需要块表达式时，在模板中显式书写 \texttt{\{ .\allowbreak{}.\allowbreak{}.\allowbreak{} yield value; \}}。

\section{参数组、变长片段及卫生绑定}

裸定义 \texttt{\$macro A} 只能用 \texttt{\$A}；空组定义 \texttt{\$macro A(\allowbreak{})\allowbreak{}} 只能用 \texttt{\$A(\allowbreak{})\allowbreak{}}。\texttt{[\allowbreak{}]\allowbreak{}}、\texttt{(\allowbreak{})\allowbreak{}}、\texttt{\{\}} 按声明顺序严格匹配，允许重复组。每组支持固定形参和一个末尾变长形参 \texttt{.\allowbreak{}.\allowbreak{}.\allowbreak{}rest} 或 \texttt{.\allowbreak{}.\allowbreak{}rest}，空捕获合法。固定形态优先于适用的变长形态。

\par
\medskip
\begin{lstlisting}[style=magel]
$macro ACCUMULATE(out)[...values]
$for value in values
int64 temporary = value;
out += temporary;
$end_for
$end_macro
int64 main() {
    int64 temporary = 100;
    int64 total = 0;
    $ACCUMULATE(total)[1, 2, 3];
    println("{} {}", total, temporary);
    return 0;
}
\end{lstlisting}
\par
\medskip

结果是 \texttt{6 100}。每轮展开中的局部声明互相独立，实参 \texttt{out} 保留调用者的绑定。变长片段可转发到另一宏，或在函数参数列表中逐项展开；\texttt{\$for} 只遍历变长形参。形参出现在声明名位置时必须传入单个标识符，该声明可供后续调用者代码使用。成员名与变量名按各自语法角色处理。

\section{宏泛型}

尖括号仅表示宏类型参数，不能作为任意文本参数组：

\par
\medskip
\begin{lstlisting}[style=magel]
$macro CONVERT<T:value>(value)
cast<T>(value)
$end_macro
$macro RESULT<T>(value)
cast<T>(value)
$end_macro
int64 main() {
    int32 source = 42;
    auto same = $CONVERT(source);
    int32 expected = $RESULT(42);
    int64 explicitValue = $RESULT<int64>(42);
    return 0;
}
\end{lstlisting}
\par
\medskip

\texttt{T:\allowbreak{}value} 只从指定实参的静态类型推导；未指定来源的 \texttt{T} 只从外层期望结果类型推导。显式类型优先，不执行参数、不自动转换、不改换推导来源。\texttt{auto x =\allowbreak{} \$RESULT(\allowbreak{}42)\allowbreak{}} 缺少期望类型，编译失败。首版不支持宏泛型值参数、特化、偏特化或任意函数的编译期执行。

\section{五类完整语法匹配}

\par
\medskip
\begin{lstlisting}[style=magel]
$macro CLASSIFY(value)
$match value is ident
1
$else is expr
2
$else is type
3
$else is block
4
$else is stmt
5
$else
0
$end_match
$end_macro
\end{lstlisting}
\par
\medskip

只支持 \texttt{expr}、\texttt{ident}、\texttt{type}、\texttt{block}、\texttt{stmt} 五类，使用同一生成解析器完整消费原始片段。按书写顺序选择首个成立分支；\texttt{ident} 也可能是 \texttt{expr}，因此顺序有意义。语法成立不代表名称存在或类型合法；只有最终展开的节点进入名称、类型、捕获、借用和资源检查。无匹配且无兜底不生成内容。未选分支的模板结构仍须合法。大括号参数组不等于块片段：传入一个完整块写 \texttt{\$A\{\{ statements; \}\}}。

\section{正则、原文与来源}

\par
\medskip
\begin{lstlisting}[style=magel]
$macro DECIMAL(value)
$regex value is "[0-9]+"
value
$else
0
$end_regex
$end_macro
\end{lstlisting}
\par
\medskip

\texttt{\$regex} 使用 C++ \texttt{std:\allowbreak{}:\allowbreak{}regex} 的 ECMAScript 语法对去除首尾空白的完整参数原文匹配，保留内部注释、引号和换行。字符串按语言规则解码一次，合法与非法模式均缓存；非法模式报编译错误。正则捕获组仅用于匹配，不导出捕获变量。

\texttt{\$stringify(\allowbreak{}value)\allowbreak{}} 返回实参原文的字符串，不先执行或展开参数。\texttt{\$concat(\allowbreak{}left,\allowbreak{} right)\allowbreak{}} 只生成单个标识符，不能生成任意代码；模板中生成的名字保持定义处上下文，不能借拼名捕获调用者局部变量。直接源码中的拼名按该源码上下文查找。\texttt{\$line}、\texttt{\$file}、\texttt{\$module} 可直接使用；模板内取调用处来源，嵌套参数转发保留原始来源。

\texttt{\$replace ADD(\allowbreak{}a,\allowbreak{} b)\allowbreak{} a + b} 是同一 AST 机制的单行定义。宏递归与规模有预算，失败诊断保留调用、定义和参数来源链。编辑器宏预览显示卫生 AST，类型提示读取最终静态类型，不显示内部绑定编号。

回归见 \href{https://gitee.com/rabbitc999/magel/blob/master/tests/magel_ast_macros.py}{宏 AST}、\href{https://gitee.com/rabbitc999/magel/blob/master/tests/magel_macro_groups.py}{分组与来源}、\href{https://gitee.com/rabbitc999/magel/blob/master/tests/magel_macro_generics.py}{静态类型} 和 \href{https://gitee.com/rabbitc999/magel/blob/master/tests/magel_macro_branches.py}{语法及正则}。

\section{函数、泛型与宏如何选择}

只需计算并返回一个值时，优先函数；算法相同而类型不同，考虑泛型；需要组合调用处语法、生成声明或重复语句时，再考虑宏。宏展开在编译期完成，不等于模板内所有函数调用都在编译期执行。

\texttt{\$SQUARE(\allowbreak{}next(\allowbreak{})\allowbreak{})\allowbreak{}} 展开后的两个操作数都含 \texttt{next(\allowbreak{})\allowbreak{}}，因而可能调用两次。如果必须只求值一次，在宏模板内部显式建立局部变量或改用函数；卫生性只保护名称绑定，并不改变求值次数。

\section{一个自己声明临时列表的 vec}

下面的宏通过首项推导 \texttt{T}，在显式块中创建临时列表，最后 \texttt{yield} 列表值。宏调用仍处在普通的类型与借用检查之下。

\par
\medskip
\begin{lstlisting}[style=magel]
$macro vec<T:first>[first,...rest]
{
    List<T> result;
    result.push(first);
    $for element in rest
    result.push(element);
    $end_for
    yield result;
}
$end_macro

int64 main()
{
    List<int64> result;
    result.push(99);
    List<int64> values = $vec[1, 2, 3];
    println("{} {} {}", values.len(), values.at(2), result.at(0));
    return 0;
}
\end{lstlisting}
\par
\medskip

\par
\medskip
\begin{lstlisting}[style=out]
3 3 99
\end{lstlisting}
\par
\medskip

调用方也有 \texttt{result}，但它和宏内部临时列表没有成为同一个变量。宏内部名称由定义与卫生绑定决定，实参则保留调用处绑定。示例至少需要一项，因为首项负责类型推导；空列表应另行设计显式类型形态，不要假定本宏已经支持。

\section{定位展开错误}

将模板内部 \texttt{List<\allowbreak{}T>\allowbreak{} result;} 删除，再在调用方声明同名变量，不会使模板里的自由名称合法。应首先按名称解析问题处理，而不能被后续 \texttt{yield} 与其他节点的类型错误掩盖。

排查时依次核对调用分隔符、固定与变长参数、类型来源、模板绑定和结果类型。宏预览用于查看展开结构，真正的合法性仍由最终语义检查确定。语法匹配成功不保证名称存在，也不保证借用可逃逸。

\subsection*{故意失败示例}

下面的模板引用未声明名称，即使 main 内存在同名变量也应拒绝。它是诊断习题，不是可运行程序。

\par
\medskip
\begin{lstlisting}[style=magel]
$macro broken(value)
{
    temp.push(value);
    yield temp;
}
$end_macro

int64 main()
{
    List<int64> temp;
    List<int64> values = $broken(1);
    return 0;
}
\end{lstlisting}
\par
\medskip

\section{边界与验证}

宏操作保留语法和来源，但不会自动缓存重复出现的实参。卫生绑定避免意外捕获，不能用调用者恰有同名变量来补救模板中的未定义名称。

\section{习题}

前两题用于理解和基本实现，第三题检查边界，最后一题为综合挑战。先独立完成，再核对参考解答。

\begin{enumerate}

\item 为什么 SQUARE(next()) 可能执行两次 next？

\item 写一个显式接收目标变量的累加宏。

\item 解释用户给出的 vec 宏中 temp 未定义的原因。

\item 挑战：把 vec 宏改为自己拥有临时列表的块表达式。

\end{enumerate}

\subsection*{参考解答与验收要点}

\begin{enumerate}

\item 模板中参数出现两次，每个替换位置都可能求值；AST 宏并不自动记忆实参结果。

\item 目标作为形参传入，循环展开逐项累加，验证空参数及多个参数。

\item 模板未声明该名称且没有对应形参，调用处的同名局部变量不会自动被捕获，应优先报告名称解析错误。

\item 在模板内部显式声明 List<T> 临时值，追加实参并 yield；检查推导来源、空列表策略和返回值类型。

\end{enumerate}

\section{番外：宏展开之后}

改编自《聊斋志异·画皮》，以下为编程主题的独立改写。

书生得到一张宏皮，展开前不过几行，能替他生成一整段程序。他高兴地把它贴到每处重复代码上，连一次有副作用的调用也填了进去。

第一次运行，计数器跳了两次。书生怪展开器暗中动手，道士却把模板揭开：参数原样放在两个位置，当然发生两次求值。漂亮的外表从没承诺替他缓存结果。

书生又在外面摆一个名叫 temp 的变量，要求宏里未声明的 temp 自动来帮忙。道士摇头，说模板有自己的绑定来处，不能看见同名就随意夺取调用者的局部对象。

于是书生重新写模板：需要调用者提供的对象就设成参数，需要临时值就自己声明。展开后的结构仍要经过名称、类型与借用检查。后来他再见几行能变出百行的宏，也先问里面究竟生成了什么，不肯只看那张皮有多光鲜。

\subsection*{技术对照}

宏操作保留语法和来源，但不会自动缓存重复出现的实参。卫生绑定避免意外捕获，不能用调用者恰有同名变量来补救模板中的未定义名称。
